% Noether P09 Korean producer draft — UNCHECKED
% T01-U04; ED0002 lines6385--6404; source 948 LF B / 0C4812C43AA17340CC1E096FBF3A496C7D72A25D71A1B597EF0A09AD7B97576B
% Pointer v009 / ED0002: B06BE3530D9CF2E82B56FDBA7FE41D5D044DF2425DFA2A059D4939EAA2F7A6C2 / C9A125167ACB33D914EE4374B65AE7CDF0052F568371B8B77B720EA178ABF0E3
% Translation only; no review, build, render, or approval.

$\mathfrak G_\eta$는 다음을 포함한다.
\begin{enumerate}
\item[1)] 모든 기저의 수 $\eta$,
\item[2)] 정수적인\srcfn{***)}{이 글에서 ``정수적''이라는 말은 언제나 계수가 모든 대수적 정수의 정역 $[K]$에 속한다는 뜻으로 이해한다. $\mathfrak G_\eta$를 정의할 때에는 1)부터 5)까지에서 유리정수 계수만을 고려해도 충분하다. 그러나 더 일반적인 영역에서는 그렇게 하면 기저 성질들이 덜 단순해진다.} 계수를 갖는 $\eta$의 모든 다항식,
\item[3)] 정수적인 계수를 갖는 $\eta$의 모든 정수적 대수함수.
\end{enumerate}
$\mathfrak G_\eta$는 다음을 배제한다.
\begin{enumerate}
\item[4)] 정수적인 계수를 갖는 $\eta$의 모든 분수형 유리함수,
\item[5)] 정수적인 계수를 갖는 $\eta$의 모든 분수형 대수함수.\srcfn{*)}{Zermelo, 앞의 글, § 3.}
\end{enumerate}
